
\begin{abstract}
%\shafiq{Let's put a project page}


%\zixuan{Title: MAS-Orchestra: Elevating Collective Intelligence via Multi-Agent System Orchestration — A Training Framework, Benchmark, and Analysis?}

% \zixuan{working on Abstract}

多智能体系统（multi-agent systems，MAS）有望通过智能体间的协作提升智能水平，但现有自动化 MAS 设计方法的实际表现仍未达到预期。这些不足源于两个关键因素：（1）方法复杂性——智能体编排采用顺序式代码级执行，既限制了全局系统层面的整体推理，也难以随智能体复杂度扩展；（2）效能不确定性——人们在尚未弄清 MAS 相较单智能体系统（single-agent systems，SAS）是否具有切实收益的情况下便部署了 MAS。我们提出 \ourframework，这是一个训练时框架，它将 MAS 编排表述为采用整体式编排的函数调用强化学习问题，一次性生成完整的 MAS。在 \ourframework 中，复杂且面向目标的子智能体被抽象为可调用函数；这种抽象在隐藏内部执行细节的同时，使模型能够对系统结构进行全局推理。为了严谨研究 MAS 在何时以及为何有益，我们提出受控基准 \ourbenchmark，从 \depth、\horizon、\breadth、\myparallel 和 \robustness 五个维度刻画任务。我们的分析表明，MAS 的收益并非普遍存在，而是关键取决于任务结构、验证协议以及编排器与子智能体双方的能力。以这些洞见为指导，\ourframework 在数学推理、多跳问答和基于搜索的问答等公共基准上取得了稳定提升，{同时效率比强基线高出超过 \textit{10$\times$}}。\ourframework 与 \ourbenchmark 共同推动了 MAS 的训练与理解，助力实现多智能体智能。
% \footnote{We will open-source the data, checkpoint, code, leaderboard for all components upon acceptance.}
% \footnote{Code, data, examples and the trained orchestrators will be publicly available upon the acceptance}




%Multi-agent systems (MAS) promise stronger collective intelligence by enabling coordination among agents. However, existing approaches to automatic MAS design often formulate orchestration as executable code that does not scale to complex agents, rely on sequential decisions that prevent global system-level reasoning, and lack principled understanding of when multi-agent coordination is beneficial over single-agent systems (SAS).



%We propose \ourframework, a training-time framework that formulates MAS orchestration as a function-calling reinforcement learning problem with holistic orchestration, generating a complete MAS in a single step. This formulation encapsulates complex, goal-oriented sub-agents as callable functions, enabling global reasoning over system structure without exposing sub-agent internals. To systematically analyze MAS benefits over SAS, we introduce \ourbenchmark, a controlled benchmark that characterizes tasks along five axes: \depth, \horizon, \breadth, \myparallel, and \robustness. Using this benchmark, we conduct rigorous analysis that reveals MAS effectiveness depends on task structure, verification protocols, and the capabilities of both the orchestrator and sub-agents, rather than holding universally. Guided by these insights, \ourframework achieves consistent improvements on public benchmarks, including mathematical reasoning, multi-hop question answering, and search-based question answering. 
%Together, \ourframework, \ourbenchmark, and comprehensive analysis provide a principled foundation for designing and understanding MAS that genuinely elevate collective intelligence.


%\austin{Abstract is a bit lengthy, 1/18 attempt to condense}

 %This work addresses both factors: We propose \ourframework, a simple and flexible training framework that formulates agent orchestration as a function-calling reinforcement learning (RL) problem. The formulation of goal-oriented sub-agents as callable functions enables global reasoning over system structure.  We also propose \ourbenchmark, a benchmark engine that enables analysis of MAS and SAS over five axes: \depth, \horizon, \breadth, \myparallel, \austin{Parallel maybe should be Parallelism (Noun)} and \robustness. Using \ourbenchmark, we rigorously analyze specific factors such as task structure, verification protocols, and orchestrator/sub-agent interplay to guide our MAS development. In all, MAS trained with \ourframework yield consistent improvements in math reasoning, multi-hop question answering, and search-based question answering. 
 
 
\end{abstract}
